\MFQTransition{事件仿真、策略比较与综合项目题组}
\begin{enumerate}[leftmargin=*]
  \item \textbf{口述概念：}为什么共同随机数能提高控制器差异比较的精度？
  \item \textbf{笔试推导：}写出静态半点差策略在给定方向与价格变化下的逐事件 PnL。
  \item \textbf{统计设计：}为 PnL、期末库存和最大库存设计置信区间。
  \item \textbf{数值编程：}把齐次 Poisson 换为日内季节强度，并验证时间变换残差。
  \item \textbf{反例实验：}重复记录一次成交或改变事件先后顺序，手算库存与现金会怎样出错，并指出被破坏的数量或信息关系。
  \item \textbf{研究判断：}审阅只报告单路径 PnL 的做市研究材料。
  \item \textbf{真实数据：}解释 SEC 聚合 cancel-to-trade 比率为何不能识别自己的队列成交概率。
  \item \textbf{综合项目：}建立报价、成交与库存相互影响的小模型，比较两种策略，并解释参数变化怎样影响结论。
\end{enumerate}

\subsection*{基础衔接：独立计算}
两策略在同一路径上的损益完全相同，配对差标准误为零能否证明市场盈利？
